\subsection{Defect Imaging Comparison}
\label{sec:experiment-imaging}

To evaluate the ability of each method to preserve defect-relevant features in full waveform images, a B-scan reconstruction experiment was conducted over the defect-containing aluminum specimens described in Section~\ref{sec:experiment-setup}. Ultrasonic signals were acquired over a regular grid spanning the specimen surface, covering the machined slots and notches. Each denoising method was applied independently to every scan location, and the resulting signals were used to reconstruct B-scan images.

\textbf{Denoising methods compared.} Signal averaging with 256 shots per location served as the ground-truth reference for B-scan reconstruction. The denoised output of each method was used to form two image types: (1) the peak-to-peak amplitude of the direct Lamb wave arrival, and (2) the Hilbert envelope of the gated waveform, sampled at the theoretical arrival time predicted from the specimen geometry. The proposed PMDF-Net was applied in a single forward pass using the trained model weights. Wiener filtering estimated the noise power spectral density from the signal autocovariance averaged over a 201-sample window. BM3D used a block-matching similarity threshold of 3.5 and a hard thresholding strategy adapted to 1D ultrasonic waveforms. DWT applied soft thresholding with the db4 wavelet at a decomposition level selected by five-fold cross-validation on the noisy training set. All parameters were held fixed during the imaging experiment.

\textbf{Evaluation protocol.} Image quality was assessed against the high-averaged reference B-scan using three complementary metrics. Full-reference metrics comprised SNR improvement ($\Delta$SNR, defined in Section~\ref{sec:metrics}) and normalized MSE (NMSE), computed pixel-wise across the full image extent. Feature preservation metrics evaluated the arrival time estimate and peak amplitude of the direct Lamb wave arrival at four manually identified locations near each defect boundary; errors are reported as mean $\pm$ standard deviation across all defect sites. Defect imaging quality was quantified using the contrast-to-noise ratio (CNR) of each defect region relative to a noise-only background region:
\begin{equation}
\text{CNR} = \frac{|\mu_{\text{defect}} - \mu_{\text{bkg}}|}{\sqrt{\sigma_{\text{defect}}^2 + \sigma_{\text{bkg}}^2}},
\end{equation}
where $\mu$ and $\sigma$ denote the mean and standard deviation of the pixel intensities in the defect and background regions, respectively. CNR was computed separately for each machined slot and notch, then averaged. Statistical significance of metric differences across methods was assessed using the Wilcoxon signed-rank test with Holm--Bonferroni correction ($p < 0.05$).
